\subsection{Analyzing Arguments with Truth Tables}

\content{
        \begin{example} Consider the statement ``If you cook salmon in the break
        room oven, I will be mad at you". What are the converse, inverse, and
        contrapositive statements?
        \end{example}
}{% presentation
\vspace{2in}
}{% nopresentation
\vspace{1.5in}
}{% comments
}


\content{
        \begin{definition} (Biconditional) The biconditional statement is used
        when two statements are interchangeable. Denoted $p\Leftrightarrow q$.

        The truth table of the biconditional is as follows: 

        \begin{tabular}{|c|c|c|}
        \hline
        $p$ & $q$ & $p\Leftrightarrow q$  \\
        \hline
        T & T & T \\
        \hline
        T & F & F \\
        \hline
        F & T & F \\
        \hline
        F & F & T \\
        \hline
        \end{tabular}

        This is logically equivalent to $(p\Rightarrow q )\land (q\Rightarrow p)$. In
        words this is often written ``$p$ if and only if $q$".
        \end{definition}
}{% presentation
}{% nopresentation
}{% comments
}

\content{
        \begin{example} Suppose that this statement is true: ``The garbage truck
        comes down my street if and only if it is Tuesday morning." Which of the
        following statements could be true? 
        \begin{enumerate}
        \item It is noon on Tuesday, and the garbage truck did not come down my
        street this morning.
        \item It is Monday, and the garbage truck came down my street this
        morning.
        \item It is Wednesday, and the garbage truck did not come down my street
        this morning.
        \end{enumerate}
        \end{example}
}{% presentation
}{% nopresentation
}{% comments
}

\content{
        \begin{definition} (Analyze Argument with Truth Table) To perform the
        analysis, we do the following steps:
        \begin{enumerate}
        \item Represent each of the premises symbolically
        \item Combine all premises to form a compound statement, this will form
        the antecedent
        \item Create a conditional statement with the conclusion as the
        consequent
        \item Create a truth table for the statement, if it is always true, the
        argument was valid.
        \end{enumerate}
        \end{definition}
}{% presentation
}{% nopresentation
}{% comments
}

\content{
        \begin{example} Consider the argument: 
        \begin{enumerate}
        \item Premise: ``If you bought bread, you went to the store."
        \item Premise: ``You bought bread"
        \item Conclusion: ``You went to the store."
        \end{enumerate}
        \end{example}
}{% presentation
\begin{tabular}{|c|c|c|c|}
\hline
 &  &\hspace{0.75cm} &\hspace{0.75cm}  \\
\hline
T & T &&  \\
\hline
T & F &&  \\
\hline
F & T &&  \\
\hline
F & F &&  \\
\hline
\end{tabular}
}{% nopresentation
\begin{tabular}{|c|c|c|c|}
\hline
 &  &\hspace{0.75cm} &\hspace{0.75cm}  \\
\hline
T & T &&  \\
\hline
T & F &&  \\
\hline
F & T &&  \\
\hline
F & F &&  \\
\hline
\end{tabular}
}{% comments
}

\content{
        \begin{example} Consider the argument: 
        \begin{enumerate}
        \item Premise: ``If you listen to the Grateful Dead, you are a hippie."
        \item Premise: ``Sky does not listen to the Grateful Dead"
        \item Conclusion: ``Sky is not a hippie."
        \end{enumerate}
        \end{example}
}{% presentation
\begin{tabular}{|c|c|c|c|}
\hline
 &  &\hspace{0.75cm} &\hspace{0.75cm}  \\
\hline
T & T &&  \\
\hline
T & F &&  \\
\hline
F & T &&  \\
\hline
F & F &&  \\
\hline
\end{tabular}
}{% nopresentation
\begin{tabular}{|c|c|c|c|}
\hline
 &  &\hspace{0.75cm} &\hspace{0.75cm}  \\
\hline
T & T &&  \\
\hline
T & F &&  \\
\hline
F & T &&  \\
\hline
F & F &&  \\
\hline
\end{tabular}
}{% comments
}

\content{
        \begin{example} Consider the argument: 
        \begin{enumerate}
        \item Premise: ``If I work hard, I will get a raise"
        \item Premise: ``If I get a raise, I will buy a boat"
        \item Premise: ``I did not buy a boat"
        \item Conclusion: ``I did not work hard"
        \end{enumerate}
        \end{example}
}{% presentation
\begin{tabular}{|c|c|c|c|c|c|c|}
\hline
p & q & r &\hspace{0.75cm} & \hspace{0.75cm} & \hspace{0.75cm} & \hspace{0.75cm}  \\
\hline
T & T & T &&&&\\
\hline
T & T & F &&&&\\
\hline
T & F & T &&&&\\
\hline
T & F & F &&&&\\
\hline
F & T & T &&&&\\
\hline
F & T & F &&&&\\
\hline
F & F & T &&&&\\
\hline
F & F & F &&&&\\
\hline
\end{tabular}
}{% nopresentation
\begin{tabular}{|c|c|c|c|c|c|c|}
\hline
p & q & r &\hspace{0.75cm} & \hspace{0.75cm} & \hspace{0.75cm} & \hspace{0.75cm}  \\
\hline
T & T & T &&&&\\
\hline
T & T & F &&&&\\
\hline
T & F & T &&&&\\
\hline
T & F & F &&&&\\
\hline
F & T & T &&&&\\
\hline
F & T & F &&&&\\
\hline
F & F & T &&&&\\
\hline
F & F & F &&&&\\
\hline
\end{tabular}
}{% comments
}
